\bookchapter{Chapitre 33 : Perspectives futures — La prochaine étape du smartphone}

\chapterintro{Le smartphone ne disparaîtra pas du jour au lendemain, mais sa « position hégémonique » pourrait être partagée avec de nouveaux appareils. Ce chapitre aborde trois sujets : à quoi ressembleront les appareils de nouvelle génération, comment le smartphone « s'effacera » pour devenir la pierre angulaire de l'identité numérique, et les questions éthiques que nous devons anticiper.}

\booksection{I. AR/VR et calcul spatial : le « second bureau » au-delà du smartphone}

La concurrence des appareils de nouvelle génération a commencé ; les plus susceptibles de capter le temps du smartphone sont les appareils AR/VR et de calcul spatial :

\begin{itemize}
  \item \textbf{Successeur du Vision Pro}: Le lancement de l'Apple Vision Pro marque l'arrivée du concept de « calcul spatial », mais le prix de la première génération (à partir de 29 999 yuans) et le poids restent des obstacles. Une version Vision plus légère et moins chère est attendue en 2026-2027, faisant passer le calcul spatial du « jouet de développeur » à « l'adoption précoce grand public ». L'iPhone pourrait alors devenir le compagnon de calcul et le centre de données des lunettes Vision.
  \item \textbf{Lunettes AR}: Ouvre un « second bureau » devant vos yeux, les informations numériques deviennent omniprésentes comme l'air. Les lunettes intelligentes Meta x Ray-Ban se sont vendues à des millions d'exemplaires, et les fabricants nationaux Rokid, XREAL, INMO lancent également des lunettes AR grand public. Les lunettes AR de 2025-2026 peuvent faire de la traduction en temps réel, de la navigation HUD et du dialogue avec un assistant IA, mais restent limitées par l'affichage à guide d'ondes et l'autonomie, à 3-5 ans de « remplacer le smartphone ».
  \item \textbf{Bagues intelligentes et wearables}: Samsung Galaxy Ring, Oura Ring et autres bagues intelligentes gagnent en popularité, servant de « capteurs corporels » du smartphone pour compléter le suivi de santé. Le rôle du smartphone passe de « terminal universel » à une combinaison « hub + accessoires légers multiples ».
\end{itemize}

\booksection{II. Interface cerveau-machine : la pensée comme moyen d'interaction}

La pensée elle-même deviendra un moyen d'interaction, mais le chemin est plus long qu'on ne l'imagine :

\begin{itemize}
  \item \textbf{Dernières avancées de Neuralink}: Le Neuralink d'Elon Musk a réalisé en 2024 sa première implantation humaine ; le premier sujet, Noland Arbaugh, contrôle une souris d'ordinateur par la pensée pour jouer aux échecs. Actuellement, les BCI (interfaces cerveau-machine) implantables sont principalement destinés à l'assistance des paralysés, dans le but de les aider à « manipuler » à nouveau le monde numérique.
  \item \textbf{BCI non invasif}: Plusieurs entreprises développent des casques non invasifs (Kernel, OpenBCI), capables de lire les signaux cérébraux sans chirurgie, mais avec une précision bien moindre que les versions implantables. « Taper par la pensée » est encore lointain.
  \item \textbf{Perspectives grand public}: L'interface cerveau-machine est à 10-20 ans de l'interaction grand public sur smartphone, mais pourrait d'abord servir à l'accessibilité et au suivi de la concentration. Les questions éthiques et de confidentialité (« à qui appartiennent mes ondes cérébrales ») seront débattues avant la maturité technologique.
\end{itemize}

\booksection{III. Le smartphone « s'efface » et devient la pierre angulaire de l'identité}

Une autre direction d'évolution du smartphone pourrait être de « s'effacer » — devenir un « facilitateur invisible » omniprésent, la unique pierre angulaire de votre identité numérique :

\begin{itemize}
  \item \textbf{Authentification à deux facteurs dans le cloud}: Authentification « vous êtes vous » instantanée, connexion à tout service sans saisie manuelle de code.
  \item \textbf{Patrimoine numérique comme réglage d'usine}: Préréglage d'un plan de succession des actifs numériques lors de l'activation d'un nouvel appareil, faisant de la « fin de vie » un paramètre par défaut.
  \item \textbf{Du « centré sur l'appareil » au « centré sur l'identité »}: Votre coffre-fort de données personnelles dans le cloud est le centre, le matériel n'est qu'un point d'accès. En changeant d'appareil, votre identité et vos données vous suivent.
\end{itemize}

\booksection{IV. Réflexion éthique : sommes-nous encore les maîtres du smartphone ?}

Une question ultime persiste : sommes-nous des « dieux » habilités par la technologie, ou des « cerveaux en cuve » nourris par des algorithmes ?

\begin{itemize}
  \item \textbf{Éthique de l'IA}: Quand l'IA côté appareil connaît vos préférences mieux que vous, vos « choix » sont-ils le fruit de votre volonté ou le résultat d'une guidance ? Quand un AI Agent passe une commande ou répond à un message pour vous, où sont la responsabilité et les limites ?
  \item \textbf{Souveraineté des données}: Vos données de santé, vos ondes cérébrales, vos données biométriques — vous appartiennent-elles, au fabricant ou au fournisseur de cloud ? La « Loi sur la protection des informations personnelles » et le RGPD européen ne sont qu'un point de départ ; le « droit à la portabilité » et le « droit à l'oubli » ont encore un long chemin avant d'être réellement appliqués.
  \item \textbf{Fracture numérique}: Les smartphones IA, les écrans pliables et la communication par satellite creusent l'écart entre ceux qui « peuvent se le permettre » et ceux qui « ne le peuvent pas ». L'accès équitable aux infrastructures numériques deviendra un nouvel enjeu social — tout aussi crucial qu'il y a 20 ans, « savoir ou non utiliser Internet ».
  \item Quand l'algorithme sait mieux que vous ce que vous voulez voir, dans quelle mesure vos « j'aime » sont-ils spontanés ?
  \item Quand vous allumez inconsciemment l'écran toutes les quelques minutes, est-ce votre habitude ou un « comportement addictif » conçu ?
\end{itemize}

\booksection{V. Rester connecté, rester indépendant — votre philosophie de vie numérique}

Quelle que soit l'évolution technologique, ces principes de vie numérique resteront valables dans 10 ans :

\begin{enumerate}
  \item Soyez distributeur de ressources, pas récepteur de notifications push.
  \item Faites de l'automatisation votre fidèle « copilote ».
  \item Construisez votre solide « arrière-garde numérique ».
  \item Pratiquez régulièrement le « désencombrement numérique » et le « jeûne social ».
  \item Considérez toujours la sécurité et la vie privée comme le préalable à toute commodité.
\end{enumerate}

\begin{tcbpractice}{Essayez : rédigez vos « Conditions d'utilisation de l'appareil numérique personnel »}
Rédigez pour vous et votre smartphone (ainsi que vos futures lunettes AR, assistant IA) de simples « conditions d'utilisation » pour définir votre relation. Par exemple :

\begin{itemize}
  \item Article 1 : Mon temps d'attention m'appartient. Tu dois passer en mode « Ne pas déranger » quand je travaille ou suis avec ma famille.
  \item Article 2 : Mes données personnelles sont ma propriété. Ta première responsabilité est le chiffrement sécurisé et la synchronisation pratique.
  \item Article 3 : Notre objectif commun est d'améliorer l'efficacité et l'expérience de vie. Si tu me rends dépendant par des designs manipulateurs, j'ai le droit de te désinstaller.
\end{itemize}

\end{tcbpractice}

\begin{tcbsummary}{Résumé du chapitre}
\begin{itemize}
  \item Le Vision Pro et les lunettes AR sont de nouveaux points d'accès au calcul au-delà du smartphone, mais ne le remplaceront pas à court terme
  \item L'interface cerveau-machine est à 10-20 ans du grand public, les questions éthiques précèdent la maturité technologique
  \item Le smartphone évolue vers « l'effacement », passant du centre matériel au centre d'identité
  \item Éthique de l'IA, souveraineté des données, fracture numérique : les enjeux centraux des 10 prochaines années
  \item Rester connecté tout en restant indépendant, la philosophie immuable de la vie numérique
\end{itemize}

\end{tcbsummary}
